\subsection{The published configuration reproduces nothing that can be compared}

Under the original's own recipe (Table~\ref{tab:blockA}) every architecture
returns the same numbers to three decimal places: \Awsdgruiou{} IoU on WSD for
the gated recurrent unit, the long short term memory, the plain recurrent
network, the 3D convolutional network and the echo state network alike. The
cause is visible in the predictions rather than the metrics: the models emit a
constant 0.5 everywhere, so thresholding produces one fixed mask for every
input. Weight decay of 0.1 applied at a learning rate of 0.1 shrinks the
encoder faster than the loss can shape it.

Two consequences follow. First, the differences between cells reported in the
original are differences between degenerate models. Second, the white noise
control (\Awsdnoiseiou{} IoU) is not a floor that a collapsed model clears by
learning; it is simply what a fixed mask scores against a slowly moving target.

The copy control settles the point. Predicting that nothing changes scores
\Awsdcopyiou{} IoU on WSD and \Absdcopyiou{} on BSD, above every figure in the
original study. Whatever those models were ranked on, it was not the ability to
predict the level set.

\subsection{The optimiser and the head are separable causes}

Replacing the optimiser alone (Table~\ref{tab:blockF}) removes the collapse:
the same fully connected head now reaches \Fwsdgruiou{} IoU for the gated
recurrent unit and \Fwsdesniou{} for the echo state network, with change IoU of
\Fwsdgruchangeiou{} and \Fwsdesnchangeiou{} respectively. The architecture was
never the reason nothing trained.

Replacing the head as well (Table~\ref{tab:blockG}) lifts every model to about
\Gwsdgruiou{} IoU. The gap between the two heads is far larger than any gap
between recurrent cells, which is the central measurement of this paper: the
original compared cells through a bottleneck that dominated the result.

\subsection{Once the task discriminates, the cells do not}

At a ten iteration horizon with a four frame window, scoring only the pixels
that move, the ordering the original reported does not appear at all
(Table~\ref{tab:blockG}). On WSD the change IoU values are
\Gwsdgruchangeiou{} (CGRU), \Gwsdlstmchangeiou{} (CLSTM),
\Gwsdrnnchangeiou{} (CRNN), \Gwsdesnchangeiou{} (CESN) and
\Gwsdlsmchangeiou{} (CLSM), a spread comparable to the seed to seed standard
deviation of each arm. On BSD the same holds. An untrained random reservoir,
analogue or spiking, matches networks trained end to end.

The original reported the echo state network at roughly a third of the gated
recurrent unit's IoU. That gap is not reproduced here under any configuration
in which both models train at all.

\subsection{What the cheap readout costs}

Fitting only a linear readout in closed form, with no backpropagation anywhere
and a frozen random encoder, takes \Gwsdesnridgefitseconds{} seconds on WSD
against \Gwsdesnfitseconds{} seconds for the gradient trained reservoir, a
reduction of more than an order of magnitude. It is also clearly worse:
\Gwsdesnridgechangeiou{} change IoU against \Gwsdesnchangeiou{}. The spiking
liquid loses more than the analogue reservoir does
(\Gwsdlsmridgechangeiou{}), which is consistent with its state being a
filtered spike train rather than a continuous activation.

This is the honest form of the reservoir argument. A reservoir is cheap when it
is fitted the way reservoirs are meant to be fitted, and on this task that
cheapness costs accuracy. Training the readout by gradient descent, as the
original did, buys the accuracy back and gives up the cost advantage: its
fitting time is then the same as a trained recurrent network's.

\subsection{Rollout separates conservative models from dynamic ones}

Under free running rollout the global and change metrics move in opposite
directions. The gated models keep a higher whole image IoU
(\Gwsdgrurolloutiou{} for CGRU) while the 3D convolutional network and the
reservoirs score higher on the moving front (\GwsdTdcnnrolloutchangeiou{}
and \Gwsdesnrolloutchangeiou{} against \Gwsdgrurolloutchangeiou{}). A model
that predicts little change tracks a slow target well in aggregate and captures
none of its motion. Reporting only one of these numbers would support either
conclusion, which is why both are in the table.

\subsection{The parameters the original blamed do not move the result}

Sweeping the spectral radius from 0.5 to 1.5 and the leaking rate from 0.0078
to 1.0, in a configuration where the model trains, moves change IoU between
\esnsweepmin{} and \esnsweepmax{} across \esnsweepcells{} cells, a standard
deviation of \esnsweepsd{} (Figure~\ref{fig:sweep}). For comparison, three
seeds of one fixed configuration vary by \seedrepsd{}. The reservoir
parameters the original identified as the cause of failure change the outcome
less than the random seed does.

The liquid's parameters behave the same way. Its firing rate responds strongly
to the threshold, from \lsmfiringmin{} to \lsmfiringmax{} spikes per neuron per
step, while its accuracy stays between \lsmsweepmin{} and \lsmsweepmax{}. The
liquid is therefore robust in accuracy and adjustable in cost, which is the
useful direction for a spiking model.

\subsection{Training on the part of the sequence that moves}

The level set is static for tens of iterations and then collapses, so uniform
sampling spends most of training on frames where nothing happens. Restricting
training to samples whose target differs from the input by at least half a
percent of pixels (Table~\ref{tab:blockH}) leaves the comparison unchanged
(CGRU \Hwsdgruchangeiou{}, CESN \Hwsdesnchangeiou{}, CLSM
\Hwsdlsmchangeiou{}) while reducing the training set. The stillness was
neither helping nor hurting; it was simply padding.
